\chapter{Design, Implementation, and Testing}
\label{ch:implementation}

This chapter presents a concise, implementation‑focused version of the design, data model, AI pipeline, clustering, visualization, and the functional test (F1) for the GJournal prototype.

\section{Scope and Objectives}
The prototype targets iOS/iPadOS/macOS and performs all NLP on device. Objectives:
\begin{itemize}
  \item Secure, local micro‑journaling with short entries.
  \item On‑device embedding and emotion detection per entry.
  \item Semantic grouping (clusters) and presentation via list, calendar, and graph.
  \item Preserve privacy: no cloud inference or storage.
  \item Maintain interactive performance on devices from iPhone 11 onward.
\end{itemize}

\section{High‑level Architecture}
GJournal uses a four‑layer architecture: \textbf{UI (SwiftUI)}, \textbf{Repository}, \textbf{Service} (AIProcessingManager, DBSCANClusterer, SafetyFilter), and \textbf{Persistence (SwiftData)}. Services are injected via protocol interfaces so mock implementations can be replaced by Core ML models without changing higher layers.

\section{Data Model}
\texttt{SwiftData} stores the primary entities:

\subsection{Bit}
Essential fields:
\begin{itemize}
  \item \texttt{id, title, feel, event, action, createdAt, lastModified}
  \item User selections: \texttt{feelEmotionRawValue, eventEmotionRawValue, actionEmotionRawValue, intensity}
  \item AI outputs: \texttt{embeddingJSON} (variable-length float array; 384 floats for MiniLM, 256 for Nomic-MRL default), \texttt{emotionProbabilitiesJSON}
  \item Flags/timestamps: \texttt{isHidden, isSafeContent, deletedAt, embeddingGeneratedAt, clusterAssignedAt}
  \item Computed accessors decode/encode JSON to typed arrays.
\end{itemize}

\subsection{Cluster}
Relational model with \texttt{@Relationship(deleteRule:.nullify)} to \texttt{Bit}. Cluster rows store \texttt{id, name, createdAt} and a list of member bits; names are updated centrally and observed reactively by SwiftData.

\subsection{Retention}
Soft delete sets \texttt{deletedAt}; a maintenance pass permanently removes entries after three days.

\section{Repository and UI Synchronization}
A repository pattern centralizes queries and writes. The repository exposes a \texttt{@Published bits: [Bit]} array; the AI pipeline sets a completion callback that reloads bits on the main thread so SwiftUI views update reactively. Search APIs accept text, emotion, date range, intensity, and cluster filters.

\section{AI Processing Pipeline}
The \texttt{AIProcessingManager} runs a six‑step pipeline on a background task (utility priority):

\begin{enumerate}
  \item \textbf{Text flattening:} combine title, Feel, Event, Action into a tagged string for the model.
  \item \textbf{Safety filter:} on‑device keyword scan across harm categories; flagged entries are stored but excluded from clustering/recommendations.
  \item \textbf{Embedding:} produce a fixed-length vector using one of three interchangeable services behind the \texttt{EmbeddingService} protocol: \texttt{CoreMLEmbeddingService} (all-MiniLM-L6-v2, 384-dim, mean pooling, 43\,MB Core ML package), \texttt{NomicEmbeddingService} (nomic-embed-text-v1.5 with MRL truncation, default 256-dim, 261\,MB Core ML package), or \texttt{NLEmbeddingService} (Apple's built-in \texttt{NLEmbedding}, 512-dim, zero bundle cost). \texttt{ServiceContainer} selects which service to wire via the \texttt{EmbeddingModelChoice} enum.
  \item \textbf{Emotion detection:} probability distribution over 12 emotion classes, produced by \texttt{CoreMLEmotionService}: MobileBERT fine-tuned on GoEmotions (28 labels collapsed to 12), converted to a 94\,MB Core ML package. Classes: Joy, Sadness, Anger, Fear, Love, Surprise, Disgust, Anxiety, Excitement, Gratitude, Hope, Guilt.
  \item \textbf{Persistence:} write embedding, emotion probabilities, safety flag, and timestamps on the MainActor.
  \item \textbf{Clustering:} re‑run DBSCAN and update \texttt{Cluster} rows.
\end{enumerate}

The pipeline minimizes UI latency by running off the main thread and switching to the MainActor only for persistence.

\section{Clustering}
\begin{itemize}
  \item \textbf{Algorithm:} DBSCAN chosen for no required cluster count and noise detection.
  \item \textbf{Distance:} cosine distance \(d(a,b)=1-\frac{a\cdot b}{\|a\|\|b\|}\).
  \item \textbf{Parameters:} adaptive epsilon calibrated per corpus via \texttt{calibrate\_epsilon} (10th percentile of pairwise distances); adaptive \texttt{minPoints} scaled by corpus size (minimum 3, maximum 11). For MRL-truncated embeddings, epsilon is further scaled by $\sqrt{d / 768}$ to compensate for the compressed distance range at lower dimensions.
  \item \textbf{Naming:} NLP-based namer extracts frequent nouns/adjectives from cluster entries; fallback to dominant emotion.
  \item \textbf{Complexity:} pairwise distance matrix \(O(N^2)\) acceptable for hundreds of entries; future optimizations (indexing/approximate NN) planned for scale.
\end{itemize}

\section{Knowledge Graph Visualization}
\begin{itemize}
  \item \textbf{Layout:} force‑directed simulation with global repulsion, intra‑cluster attraction (hub‑and‑spoke), and cluster gravity.
  \item \textbf{Initialization:} cluster‑aware sector placement to speed convergence.
  \item \textbf{Rendering:} hub‑and‑spoke edges to reduce edge count to \(O(N)\) per cluster.
  \item \textbf{Interaction:} pan, pinch‑to‑zoom, draggable nodes, and detail cards; gesture layer attached to a fixed clear background to avoid dead zones.
\end{itemize}

\section{User Interface}
Key screens and features:
\begin{itemize}
  \item \textbf{Dashboard:} greeting, activity heatmap, top emotions, quick entry.
  \item \textbf{Editor:} Feel–Event–Action fields, emotion wheel, intensity picker, auto‑save drafts.
  \item \textbf{List and Calendar:} grouped entries, full‑text search, emotion and cluster filter chips, sorting.
  \item \textbf{Graph:} interactive cluster graph with centroid labels and node detail cards.
\end{itemize}

\section{Dependency Injection and Testability}
All AI services are protocol‑based and provided by a \texttt{ServiceContainer} with three configurations: \texttt{makeMock()} (development), mixed device testing, and \texttt{makeProduction()} (Core ML models). This enables unit testing and device testing without changing application logic.

\section{Functional Test F1}
\textbf{Scope:} FR2 (create entry), FR4 (on‑device analysis), FR5 (clustering).  
\textbf{Procedure:} create a Bit in the editor, save, observe console logs, and verify calendar, list filters, and graph.  
\textbf{Expected / Observed:}
\begin{itemize}
  \item Entry persisted immediately and visible in calendar.
  \item Pipeline logs show steps 1–6; mock embedding generated (200\,ms delay).
  \item Clustering completed within seconds for small corpus; cluster name generated and filter chip appears.
  \item Knowledge graph shows the new node after clustering.
\end{itemize}

\section{Limitations and Next Steps}
\begin{itemize}
  \item Schema migration: SwiftData currently wipes the store on model changes during development. A versioned migration policy is needed before production deployment.
  \item Accessibility: VoiceOver labels on graph nodes and emotion-wheel controls require a dedicated pass.
  \item Notification rules: the scheduler logic for daily journaling reminders is defined but the exact trigger conditions (time window, streak rules) are not yet finalised.
  \item Keychain security: embedding vectors and emotion data are stored in SwiftData; a Keychain-backed encryption layer for the database file is planned.
\end{itemize}

\section{Acceptance Criteria (for this chapter)}
The implementation described here is acceptable when:
\begin{itemize}
  \item F1 is demonstrable: entry creation triggers the full pipeline and updates UI views.
  \item Data remains local and encrypted; biometric/PIN gating functions in device builds.
  \item Clustering and naming produce coherent groups on a sample dataset.
  \item Performance targets (interactive latency, incremental clustering) are met or have a documented mitigation plan.
\end{itemize}

\section{Summary}
This chapter presents the design, data model, AI pipeline, clustering algorithm, and functional test for the GJournal prototype. The system runs three interchangeable embedding services and a 12-class emotion classifier entirely on device. Remaining work is operational rather than architectural: schema migration, accessibility, and notification rules.
